% Generated draft: requires mathematical and source review.
\documentclass[11pt]{article}
\usepackage[margin=25mm]{geometry}
\usepackage{fontspec}
\setmainfont{Latin Modern Roman}
\usepackage{amsmath,amssymb,mathtools,unicode-math}
\setmathfont{Latin Modern Math}
\usepackage{enumitem,needspace,xurl,hyperref}
\setlength{\parindent}{0pt}
\setlength{\parskip}{5pt}
\setlength{\emergencystretch}{3em}
\newcommand{\N}{\mathbb{N}}
\newcommand{\Z}{\mathbb{Z}}
\newcommand{\Q}{\mathbb{Q}}
\newcommand{\R}{\mathbb{R}}
\newcommand{\C}{\mathbb{C}}
\newcommand{\F}{\mathbb{F}}
\newcommand{\A}{\mathbb{A}}
\newcommand{\PP}{\mathbb{P}}
\newcommand{\E}{\mathbb{E}}
\newcommand{\eps}{\varepsilon}
\newcommand{\dd}{\,\mathrm{d}}
\begin{document}
\section*{UnsolvedMath ID: 10600025}
\textbf{Source problem code:} catalogue entry
\subsection{Definitions and mathematical statement}
Rotational virtual knot theory uses virtual link diagrams and their permitted moves except the first virtual Reidemeister move: a virtual curl cannot be added or removed. The source reports that quantum link invariants originally defined on classical links extend to this setting, and asks for further exploration of the resulting theory.

Its described bracket construction smooths classical crossings to obtain immersed planar state curves. For such a regular oriented curve \(C\), its Whitney degree \(w(C)\) is the degree of its unit tangent map \(S^1\to S^1\); \(|w(C)|\) is unchanged by reversing orientation. The rotational bracket assigns variables according to these absolute degrees of state curves. The local weights, variable relations and normalization are not supplied. The source proposes an Arrow-Polynomial-style generalization and asks how quantum invariants behave on rotational virtual links. The mentioned example figure and cited constructions are absent.

A specifically proposed class comes from finite-dimensional quasitriangular Hopf algebras, and from earlier work with Radford using Hopf-algebra integrals. Let \(H\) be a finite-dimensional Hopf algebra over a field \(k\), with coproduct \(\Delta\), counit \(\varepsilon\) and unit \(1\). To make the integral terminology explicit, a right integral element under one conventional usage is \(\Lambda\in H\) satisfying
\[\Lambda h=\varepsilon(h)\Lambda\quad(h\in H).\]
A right integral \emph{functional} under the convolution convention \(\lambda*f=f(1)\lambda\) is \(\lambda\in H^*\) satisfying
\[(\lambda\otimes\operatorname{id})\Delta(h)=\lambda(h)1\quad(h\in H).\]
The excerpt does not say which of these dual notions or handedness conventions its categorification request uses. It asks whether finite-dimensional Hopf algebras can be categorified and, in particular, what it means to categorify a right integral. A precise answer must specify a categorical structure and decategorification recovering the selected algebra and integral identity, together with whatever compatibility is intended for the rotational-link invariants. No categorical model or required homology theory is prescribed.

The source describes the integral-based invariants as immediately extending to rotational virtual knots and deserving study in their own right. This is retained as its report; ribbon/pivotal structure, framing/normalization, field and any nondegeneracy assumptions needed in an actual construction are not supplied. The broader exploration, Arrow-style generalization, Hopf-integral study and categorification request are distinct parts of the programme.
\subsection{Short English statement}
Explore virtual knot theory with virtual curls forbidden, its Whitney-degree bracket and Arrow-style extensions, and quantum invariants from Hopf algebras. A separate target is to define a meaningful categorification of a finite-dimensional Hopf algebra and its right integral.
\subsection{Sources}
\begin{itemize}
\item Local UnsolvedMath registry, permanent ID \(10600025\), exact source code \texttt{catalogue entry}.
\item \emph{Unsolved Problems in Virtual Knot Theory and Combinatorial Knot Theory}, revised survey, problem-list item 25: \url{https://arxiv.org/abs/1409.2823}. Registry access 2026-07-29; source-tex extraction. Embedded references and the example figure are absent; the source mentions earlier work with Radford.
\item Supplied research leads: \emph{Rotational Virtual Knots and Quantum Link Invariants}, attributed in the summary to Kauffman--Manturov (2015; ``arXiv:1509.xxxx''); and \emph{Centrality and the KRH Invariant} (2021; ``arXiv:2107.xxxx''). These identifiers are placeholders, not verified citations; authors for the second lead are not supplied.
\end{itemize}
\end{document}
